\documentclass[11pt]{article}

\usepackage[a4paper,margin=1in]{geometry}
\usepackage{booktabs}
\usepackage{hyperref}
\usepackage{float}

\title{Interpretable Source and Transfer Learning with CNN, ViT, and DHVT}
\author{}
\date{March 2026}

\begin{document}

\maketitle

\begin{abstract}
This report summarizes the current saved results of the project. The study begins with a controlled CIFAR-10 comparison between a lightweight CNN, a lightweight ViT, and a compact DHVT baseline, then transfers all three architectures to EuroSAT and Brain Tumor MRI under scratch training, linear probing, and full fine-tuning. The main findings are direct: DHVT is the strongest clean CIFAR-10 classifier, the CNN remains competitive in the low-data regime, and the vanilla ViT is the most robust to texture corruption. In downstream transfer, linear probing is consistently weaker than full fine-tuning. On EuroSAT, strong scratch training is already sufficient and DHVT scratch is best. On Brain Tumor MRI, CIFAR initialization helps substantially, and pretrained DHVT with full fine-tuning is the strongest downstream model.
\end{abstract}

\section{Introduction}
CNNs use locality and weight sharing as inductive biases \cite{lecun1998gradient}. ViTs instead model patch interactions with self-attention \cite{dosovitskiy2021image}. Prior work shows that ViTs usually require more data or stronger training recipes \cite{dosovitskiy2021image,touvron2021training}, while standard CNNs often rely more on local texture cues \cite{geirhos2019imagenet}. DHVT was proposed to bridge this gap on small datasets by combining transformer-style token mixing with stronger convolutional inductive bias \cite{lu2022bridging}. This project studies these differences in two stages: first through a controlled CIFAR-10 source experiment with robustness and interpretability analysis, and then through downstream transfer to EuroSAT and Brain Tumor MRI.

\section{Experimental Setup}
\subsection{Source stage: CIFAR-10}
CIFAR-10 uses RGB images of size $32 \times 32 \times 3$ \cite{krizhevsky2009learning}. The original 50{,}000-image training split is divided into 45{,}000 training images and 5{,}000 validation images. The 10{,}000-image test split is kept fixed.

Training uses:
\begin{itemize}
    \item \texttt{RandomCrop(32, padding=4)}
    \item \texttt{RandomHorizontalFlip()}
    \item \texttt{ToTensor()}
    \item normalization with CIFAR-10 statistics
\end{itemize}

Here, \texttt{ToTensor()} and normalization are preprocessing transforms. The random crop and horizontal flip are augmentation transforms. Augmentation is applied online during loading; it does not increase the stored dataset size.

The source-stage comparison includes:
\begin{itemize}
    \item clean test accuracy
    \item data efficiency at 10\%, 25\%, 50\%, and 100\% of the training pool
    \item robustness to occlusion and texture modification
\end{itemize}

\subsection{Downstream transfer}
Two downstream datasets are used.

\textbf{EuroSAT} uses RGB images resized to $64 \times 64 \times 3$ with a stratified train / validation / test split. Horizontal and vertical flips are used as augmentation.

\textbf{Brain Tumor MRI} uses a provided \texttt{Training/} and \texttt{Testing/} folder structure. A validation split is carved from the training folder. Images are converted to RGB and resized to $128 \times 128 \times 3$. Light rotation and affine perturbation are used as augmentation.

For both downstream datasets, two modes are compared:
\begin{itemize}
    \item \textbf{scratch}: random initialization
    \item \textbf{pretrained}: initialize from the saved CIFAR-10 checkpoint, replace the classifier head, and fine-tune end to end
\end{itemize}

So in this project, \textit{pretrained} means transfer learning through CIFAR-10 initialization, not frozen-backbone linear probing.

When linear probing is used, the backbone remains frozen during training and only the downstream classifier head is updated. Gradient-based interpretability methods may still compute gradients during the visualization pass, but this is an inference-time explanation step and does not update the frozen backbone weights.

All reported source and downstream results in this document come from 100-epoch runs.

\subsection{Architecture attribution}
The repository models are lightweight research baselines, not exact reproductions of one published architecture. The CNN is a custom compact model inspired by LeNet-style convolution / pooling classifiers \cite{lecun1998gradient}, VGG-style stacked small convolutions \cite{simonyan2015very}, and Batch Normalization \cite{ioffe2015batch}. The ViT is a compact implementation inspired by the Transformer encoder \cite{vaswani2017attention} and the original Vision Transformer design \cite{dosovitskiy2021image}. The DHVT implementation is a compact reimplementation inspired by Lu et al. \cite{lu2022bridging}, using convolutional patch embedding, head-token interaction attention, and a dynamic feed-forward block. This wording is important for attribution: the code is custom, but the architectural ideas are grounded in these standard references.

\subsection{Metrics}
The project uses standard multiclass classification metrics.
\begin{itemize}
    \item \textbf{Accuracy}: fraction of correct predictions over the evaluation split.
    \item \textbf{Precision}: for one class, among predicted samples of that class, how many are correct.
    \item \textbf{Recall}: for one class, among true samples of that class, how many are correctly recovered.
    \item \textbf{F1-score}: harmonic mean of precision and recall for one class.
    \item \textbf{Macro-F1}: arithmetic mean of the class-wise F1 scores, so every class has equal weight.
    \item \textbf{Weighted-F1}: average of the class-wise F1 scores weighted by class support.
    \item \textbf{Support}: number of true samples belonging to a class in the evaluation split.
\end{itemize}

For the downstream datasets, macro-F1 is reported together with accuracy because it is more informative than accuracy alone when class balance matters. In the current downstream runs, class supports are close to balanced, so weighted-F1 is almost identical to macro-F1 and is omitted from the main tables.

\section{Results}
\subsection{CIFAR-10 source stage}
The source-stage result is now three-way. DHVT is the strongest clean classifier, while the vanilla ViT remains the most robust to texture corruption.

\begin{table}[H]
\centering
\caption{Full-data source-stage results on CIFAR-10.}
\begin{tabular}{lccccc}
\toprule
Model & Acc. & Macro Prec. & Macro Recall & Macro-F1 & Time \\
\midrule
DHVT & 88.88\% & 0.8906 & 0.8888 & 0.8891 & 59m 56s \\
CNN & 85.70\% & 0.8611 & 0.8570 & 0.8580 & 13m 56s \\
ViT & 73.36\% & 0.7330 & 0.7336 & 0.7330 & 15m 46s \\
\bottomrule
\end{tabular}
\end{table}

\begin{table}[H]
\centering
\caption{CIFAR-10 robustness on the full-data models.}
\begin{tabular}{lccc}
\toprule
Model & Clean Acc. & Occluded Acc. & Texture Acc. \\
\midrule
DHVT & 88.88\% & 83.78\% & 41.96\% \\
CNN & 85.70\% & 80.32\% & 32.98\% \\
ViT & 73.36\% & 68.81\% & 50.26\% \\
\bottomrule
\end{tabular}
\end{table}

Across data budgets, DHVT overtakes the CNN once enough data is available, while the CNN remains ahead of the vanilla ViT throughout:
\begin{itemize}
    \item DHVT: 64.42\%, 77.80\%, 82.95\%, 88.88\%
    \item CNN: 65.70\%, 75.08\%, 81.48\%, 85.70\%
    \item ViT: 53.53\%, 60.27\%, 67.03\%, 73.36\%
\end{itemize}

This refines the earlier story. The CNN is still more data-efficient than the vanilla ViT in this small-image, from-scratch setting \cite{dosovitskiy2021image,touvron2021training}. DHVT narrows that gap and becomes the strongest clean classifier at 100\% data, which is qualitatively consistent with its design goal of improving transformer behavior on small datasets \cite{lu2022bridging}. At the same time, the vanilla ViT loses much less accuracy under texture modification, which remains consistent with the texture-bias literature on CNNs \cite{geirhos2019imagenet}.

Checkpoint-based evaluation also shows that the hardest CIFAR-10 classes are \textit{cat}, \textit{dog}, and \textit{bird}, with strong \textit{cat}/\textit{dog} confusion for both models.
The DHVT rollout maps and head-token influence maps also largely overlap on CIFAR-10, which suggests that the added head-token mechanism reinforces the same discriminative regions rather than redirecting attention to qualitatively different evidence.

\subsection{EuroSAT transfer}
Transfer to EuroSAT gives only a small gain over scratch for both models.

\begin{table}[H]
\centering
\caption{EuroSAT downstream results.}
\begin{tabular}{lcccc}
\toprule
Model & Setup & Acc. & Macro-F1 & Time \\
\midrule
CNN & scratch & 96.52\% & 0.9642 & 11m 19s \\
CNN & pretrained + linear probe & 85.78\% & 0.8496 & 11m 21s \\
CNN & pretrained + full fine-tune & 96.74\% & 0.9660 & 11m 23s \\
DHVT & scratch & 97.52\% & 0.9745 & 1h 53m 56s \\
DHVT & pretrained + linear probe & 87.52\% & 0.8682 & 1h 10m 27s \\
DHVT & pretrained + full fine-tune & 96.96\% & 0.9685 & 1h 54m 19s \\
ViT & scratch & 93.59\% & 0.9331 & 31m 29s \\
ViT & pretrained + linear probe & 79.00\% & 0.7784 & 11m 43s \\
ViT & pretrained + full fine-tune & 94.15\% & 0.9385 & 31m 28s \\
\bottomrule
\end{tabular}
\end{table}

The EuroSAT result now has a clearer interpretation:
\begin{itemize}
    \item DHVT scratch is the strongest EuroSAT result at 97.52\%
    \item CNN changes very little on EuroSAT: scratch 96.52\%, linear probe 85.78\%, and full fine-tuning 96.74\%, all with roughly 11-minute runs
    \item ViT shows a stronger adaptation gap on EuroSAT: scratch 93.59\%, linear probe 79.00\%, and full fine-tuning 94.15\%
    \item CIFAR pretraining does not help DHVT on EuroSAT: DHVT full fine-tuning is slightly below its scratch result
    \item linear probing is much weaker than full fine-tuning for all transformer-style models
    \item for DHVT, linear probing saves about 44 minutes relative to full fine-tuning, but loses about 9.4 accuracy points
    \item the ViT is especially weak under a frozen backbone: 79.00\% with linear probing versus 94.15\% with full fine-tuning
\end{itemize}

This suggests that CIFAR features alone are not enough for strong EuroSAT transfer; downstream adaptation is important. At the same time, EuroSAT is easy enough that scratch training is already strong, and for DHVT the scratch run is actually the best result.

Checkpoint-based evaluation shows that the most difficult EuroSAT classes are \textit{PermanentCrop}, \textit{Highway}, and \textit{River}. The dominant confusions are \textit{PermanentCrop} $\rightarrow$ \textit{HerbaceousVegetation} and \textit{Highway} $\leftrightarrow$ \textit{River}.

\subsection{Brain Tumor MRI transfer}
The Brain Tumor MRI transfer result is stronger and more informative. Here, CIFAR pretraining gives a clear boost to both architectures.

\begin{table}[H]
\centering
\caption{Brain Tumor MRI downstream results.}
\begin{tabular}{lcccc}
\toprule
Model & Setup & Acc. & Macro-F1 & Time \\
\midrule
CNN & scratch & 82.56\% & 0.8230 & 13m 18s \\
CNN & pretrained + linear probe & 67.31\% & 0.6715 & 42m 10s \\
CNN & pretrained + full fine-tune & 89.38\% & 0.8917 & 13m 16s \\
DHVT & scratch & 87.88\% & 0.8776 & 2h 57m 00s \\
DHVT & pretrained + linear probe & 76.88\% & 0.7648 & 59m 24s \\
DHVT & pretrained + full fine-tune & 94.00\% & 0.9386 & 2h 31m 18s \\
ViT & scratch & 85.88\% & 0.8554 & 47m 47s \\
ViT & pretrained + linear probe & 76.00\% & 0.7524 & 15m 28s \\
ViT & pretrained + full fine-tune & 93.00\% & 0.9279 & 47m 56s \\
\bottomrule
\end{tabular}
\end{table}

The Brain Tumor MRI result is stronger:
\begin{itemize}
    \item DHVT full fine-tuning is now the strongest downstream result in the repository at 94.00\%
    \item CNN also benefits strongly from full fine-tuning: 82.56\% from scratch, 67.31\% with linear probing, and 89.38\% with full fine-tuning
    \item ViT shows the same pattern: 85.88\% from scratch, 76.00\% with linear probing, and 93.00\% with full fine-tuning
    \item linear probing is much weaker than both scratch and full fine-tuning
    \item for CNN, linear probing is not even cheaper than full fine-tuning here, since the frozen-backbone run still takes 42 minutes versus 13 minutes for full fine-tuning
    \item for ViT, linear probing is cheaper, but still loses about 17 accuracy points relative to full fine-tuning
    \item for DHVT, linear probing cuts training time from 2h 31m to 59m, but loses 17.1 accuracy points
    \item this drop is severe for both ViT-style models, which indicates that downstream adaptation is necessary rather than optional
    \item full fine-tuning gives a large gain over scratch
\end{itemize}

This means the project now has a split outcome:
\begin{itemize}
    \item DHVT is strongest on the source task and on EuroSAT
    \item pretrained DHVT is strongest on Brain Tumor MRI
\end{itemize}

\section{Conclusion}
The current results support three direct conclusions.

\begin{itemize}
    \item On CIFAR-10, DHVT is the best clean classifier, while the CNN remains stronger than the vanilla ViT in the lowest-data regime.
    \item On CIFAR-10, the vanilla ViT is much less affected by the texture-modified test set.
    \item For downstream transfer, linear probing is not enough. Full fine-tuning is clearly stronger on both EuroSAT and Brain Tumor MRI.
    \item CIFAR pretraining matters much more for Brain Tumor MRI than for EuroSAT, and the pretrained DHVT with full fine-tuning is currently the best medical downstream model in the repository.
\end{itemize}

So the present evidence does not support one universal winner. Instead, the results suggest that CNNs, ViTs, and DHVT learn different useful biases, and those differences matter differently across source and downstream tasks. The timing results also matter: linear probing can be cheaper for the larger transformer-style models, but the accuracy loss is too large to make it the preferred setting in the current experiments.

\begin{thebibliography}{9}

\bibitem{lecun1998gradient}
Yann LeCun, Léon Bottou, Yoshua Bengio, and Patrick Haffner.
\newblock Gradient-based learning applied to document recognition.
\newblock \emph{Proceedings of the IEEE}, 86(11):2278--2324, 1998.

\bibitem{krizhevsky2009learning}
Alex Krizhevsky and Geoffrey Hinton.
\newblock Learning multiple layers of features from tiny images.
\newblock Technical report, University of Toronto, 2009.

\bibitem{dosovitskiy2021image}
Alexey Dosovitskiy, Lucas Beyer, Alexander Kolesnikov, Dirk Weissenborn, Xiaohua Zhai, Thomas Unterthiner, Mostafa Dehghani, Matthias Minderer, Georg Heigold, Sylvain Gelly, Jakob Uszkoreit, and Neil Houlsby.
\newblock An image is worth 16x16 words: Transformers for image recognition at scale.
\newblock In \emph{International Conference on Learning Representations}, 2021.

\bibitem{simonyan2015very}
Karen Simonyan and Andrew Zisserman.
\newblock Very deep convolutional networks for large-scale image recognition.
\newblock In \emph{International Conference on Learning Representations}, 2015.

\bibitem{ioffe2015batch}
Sergey Ioffe and Christian Szegedy.
\newblock Batch normalization: Accelerating deep network training by reducing internal covariate shift.
\newblock In \emph{International Conference on Machine Learning}, pages 448--456, 2015.

\bibitem{vaswani2017attention}
Ashish Vaswani, Noam Shazeer, Niki Parmar, Jakob Uszkoreit, Llion Jones, Aidan N. Gomez, {\L}ukasz Kaiser, and Illia Polosukhin.
\newblock Attention is all you need.
\newblock In \emph{Advances in Neural Information Processing Systems}, volume 30, 2017.

\bibitem{touvron2021training}
Hugo Touvron, Matthieu Cord, Matthijs Douze, Francisco Massa, Alexandre Sablayrolles, and Hervé Jégou.
\newblock Training data-efficient image transformers \& distillation through attention.
\newblock In \emph{International Conference on Machine Learning}, pages 10347--10357, 2021.

\bibitem{geirhos2019imagenet}
Robert Geirhos, Patricia Rubisch, Claudio Michaelis, Matthias Bethge, Felix A. Wichmann, and Wieland Brendel.
\newblock ImageNet-trained CNNs are biased towards texture; increasing shape bias improves accuracy and robustness.
\newblock In \emph{International Conference on Learning Representations}, 2019.

\bibitem{lu2022bridging}
Zhiying Lu, Hongtao Xie, Chuanbin Liu, and Yongdong Zhang.
\newblock Bridging the gap between vision transformers and convolutional neural networks on small datasets.
\newblock In \emph{Advances in Neural Information Processing Systems}, 2022.

\end{thebibliography}

\end{document}
